\begin{remark}[``Complete'' is relative to a dialect]\label{rem:complete-dialect}
One qualification on ``complete''. The search is complete for the regular expressions its
parser accepts, which is the dialect of the equivalence library rather than all of
PromQL's: a pattern carrying an interior anchor, or one whose escapes the structural
proposer normalises first, is refused by the search where the proposer answers. No corpus
pattern is of either shape --- all $184$ were reached --- so the agreement reported above
is over the whole population, but the word means complete relative to that dialect.
\end{remark}
